\section*{Capítulo \ref{ch:g9:inside-the-atom} --- El interior del átomo}

\begin{solution}{exo:g9:inside-the-atom:1}
En el \omterm{def:g9:inside-the-atom:nucleus}{núcleo}: \omterm{def:g9:inside-the-atom:proton}{protones} (positivos) y \omterm{def:g9:inside-the-atom:proton}{neutrones} (sin carga). A su
alrededor: \omterm{def:g9:inside-the-atom:nucleus}{electrones} (negativos).
\end{solution}

\begin{solution}{exo:g9:inside-the-atom:2}
\ce{^{16}_{8}O}: 8 \omterm{def:g9:inside-the-atom:proton}{protones}, 8 \omterm{def:g9:inside-the-atom:proton}{neutrones}, 8 \omterm{def:g9:inside-the-atom:nucleus}{electrones}.
\ce{^{27}_{13}Al}: 13, 14, 13. \ce{^{40}_{20}Ca}: 20, 20, 20.
\end{solution}

\begin{solution}{exo:g9:inside-the-atom:3}
Tiene tantos \omterm{def:g9:inside-the-atom:nucleus}{electrones} (de carga $-e$ cada uno) como \omterm{def:g9:inside-the-atom:proton}{protones} (de carga
$+e$ cada uno): las cargas se anulan.
\end{solution}

\begin{solution}{exo:g9:inside-the-atom:4}
$A = 26 + 30 = 56$: \ce{^{56}_{26}Fe}.
\end{solution}

\begin{solution}{exo:g9:inside-the-atom:5}
Sí: los dos tienen $Z = 17$, así que los dos son cloro. Tienen distinto
\omterm{def:g9:inside-the-atom:atomic-number}{número másico} (18 y 20 \omterm{def:g9:inside-the-atom:proton}{neutrones}): son \omterm{def:g9:inside-the-atom:isotope}{isótopos}.
\end{solution}

\begin{solution}{exo:g9:inside-the-atom:6}
No: tienen distinto \omterm{def:g9:inside-the-atom:atomic-number}{número atómico} (6 y 7), así que son elementos
distintos, el carbono y el nitrógeno. Los \omterm{def:g9:inside-the-atom:isotope}{isótopos} tienen el mismo $Z$.
\end{solution}

\begin{solution}{exo:g9:inside-the-atom:7}
$56 \times 1.67 \times 10^{-27} \approx \qty{9.35e-26}{kg}$.
\end{solution}

\begin{solution}{exo:g9:inside-the-atom:8}
Un \omterm{def:g7:atoms-and-molecules:atom}{átomo} es sobre todo espacio vacío: el \omterm{def:g9:inside-the-atom:nucleus}{núcleo} es diminuto, así que la
mayoría de las partículas no encuentran ningún \omterm{def:g9:inside-the-atom:nucleus}{núcleo} en su camino y
pasan en línea recta.
\end{solution}

\begin{solution}{exo:g9:inside-the-atom:9}
$10^{-10} / 10^{-15} = 10^{5}$: cien mil veces más pequeño.
\end{solution}

\begin{solution}{exo:g9:inside-the-atom:10}
La química depende de los \omterm{def:g9:inside-the-atom:nucleus}{electrones}, y los \omterm{def:g9:inside-the-atom:isotope}{isótopos} de un elemento
tienen el mismo número de \omterm{def:g9:inside-the-atom:nucleus}{electrones} (el mismo $Z$).
\end{solution}

\begin{solution}{exo:g9:inside-the-atom:11}
$26 \times 9.11 \times 10^{-31} = \qty{2.37e-29}{kg}$. Fracción:
$2.37 \times 10^{-29} / 9.35 \times 10^{-26} \approx 2.5 \times
10^{-4}$, alrededor del \qty{0.025}{\percent} de la masa.
\end{solution}

\begin{solution}{exo:g9:inside-the-atom:12}
6915 \omterm{def:g7:atoms-and-molecules:atom}{átomos} de cobre~63 y 3085 \omterm{def:g7:atoms-and-molecules:atom}{átomos} de cobre~65. Media:
$(6915 \times 63 + 3085 \times 65)/10\,000 = 63.6$ \omterm{def:g9:inside-the-atom:proton}{nucleones}.
\end{solution}

\begin{solution}{pb:g9:inside-the-atom:1}
\textbf{1.} \ce{^{35}_{17}Cl} y \ce{^{37}_{17}Cl}.

\textbf{2.} Cloro~35: 17 \omterm{def:g9:inside-the-atom:proton}{protones}, 18 \omterm{def:g9:inside-the-atom:proton}{neutrones}, 17 \omterm{def:g9:inside-the-atom:nucleus}{electrones}.
Cloro~37: 17 \omterm{def:g9:inside-the-atom:proton}{protones}, 20 \omterm{def:g9:inside-the-atom:proton}{neutrones}, 17 \omterm{def:g9:inside-the-atom:nucleus}{electrones}.

\textbf{3.} Los dos tienen 17 \omterm{def:g9:inside-the-atom:proton}{protones}, así que los dos son el elemento
cloro; solo se diferencian en el número de \omterm{def:g9:inside-the-atom:proton}{neutrones}, así que son
\omterm{def:g9:inside-the-atom:isotope}{isótopos}.

\textbf{4.} Sí: la química depende de los \omterm{def:g9:inside-the-atom:nucleus}{electrones}, y los dos tienen
17.

\textbf{5.} $35 \times 1.67 \times 10^{-27} = \qty{5.845e-26}{kg}$.

\textbf{6.} $17 \times 9.11 \times 10^{-31} = \qty{1.55e-29}{kg}$:
unas 4000 veces menos que el \omterm{def:g9:inside-the-atom:nucleus}{núcleo}, despreciable.

\textbf{7.} $37 \times 1.67 \times 10^{-27} = \qty{6.179e-26}{kg}$.

\textbf{8.} $\qty{1}{g} = \qty{e-3}{kg}$; $10^{-3} / 5.845 \times
10^{-26} \approx 1.71 \times 10^{22}$ \omterm{def:g7:atoms-and-molecules:atom}{átomos}.

\textbf{9.} \qty{75.8}{\percent} y \qty{24.2}{\percent}.

\textbf{10.} $758 \times 5.845 \times 10^{-26} = \qty{4.430e-23}{kg}$;
$242 \times 6.179 \times 10^{-26} = \qty{1.495e-23}{kg}$.

\textbf{11.} $4.430 \times 10^{-23} + 1.495 \times 10^{-23} =
\qty{5.925e-23}{kg}$.

\textbf{12.} $5.925 \times 10^{-23} / 1000 \approx
\qty{5.93e-26}{kg}$: la masa media de un \omterm{def:g7:atoms-and-molecules:atom}{átomo} de cloro.
\end{solution}
